\section{Introduction}
\label{sec:intro}

Detecting whether a speaker has a cold from short audio clips is the kind of paralinguistic task where modern foundation-model representations should help and where the standard practice for fine-tuning them creates two persistent problems. The first is the speaker shortcut: in small-data corpora with a few hundred speakers and a binary label that correlates partially with voice identity, a classifier trained end-to-end on a foundation-model representation can score well on a per-utterance speaker recognition probe and only modestly on the actual paralinguistic task~\cite{huckvale2017cold,schuller2017interspeech}. The second is the lack of a clean methodology for measuring whether a proposed de-confounding intervention has done anything: speaker probes vary across substrates, dimensionality, and probe architecture, and an apparent ``probe top-1 dropped from $X$ to $Y$'' result can be an artefact of the bottleneck the intervention introduced rather than de-confounding signal.

This paper documents a self-contained pipeline on the URTIC corpus from the ComParE 2017 Cold sub-challenge~\cite{schuller2017interspeech} that addresses both problems while accepting both the architectural commitment of a frozen WavLM-Large~\cite{chen2022wavlm} backbone and a small ($\sim$8.5\,k chunks, $k{=}210$ pseudo-speakers) compute budget. The pipeline consists of (i) a layer-weighted classification head trained on cached pooled statistics of the frozen backbone (\textsc{A2}); (ii) an audit step that scores each WavLM layer on cold-vs-speaker decoupling (\textsc{A5d}); (iii) a re-initialisation of \textsc{A2}'s layer-weight softmax with the per-layer audit signal as a prior, producing a strict improvement on UAR with no probe inflation (\textsc{A2.5}, $+0.020$ UAR $\sim$7$\sigma$); (iv) an honesty-audited late-fusion stage that admits low-dimensional handcrafted feature groups by a subtractive cold-vs-speaker score and combines them with the WavLM head's logit at a fixed $\beta$ (\textsc{A5b}); and (v) a categorical evaluation of three classes of de-confounding intervention (data-level cross-speaker embedding mixup, representation-level supervised contrastive with speaker-masked positives, gradient-level DANN/MDD adversary) at every architectural scope tractable in the project's compute budget. The de-confounding ladder closes negatively at every scope: the speaker shortcut on URTIC + frozen WavLM is mechanism-resistant at the strengths tested.

\paragraph{Contributions.} We argue three contributions.

\textbf{C1: A controlled positive system.} \textsc{A2.5 + A5b K=2 (G4\_gain\_invariant + G5\_modulation)} reaches devel\_test UAR $0.7037 \pm 0.0060$ at $N{=}5$, a $+0.066$ cumulative lift over a leak-corrected uniform-A2 baseline ($\sim$17$\sigma$ at $N{=}5$). A 5-seed mean-logit ensemble pushes this to $0.7090$, within $0.001$ UAR of the original 2017 baseline of $0.710$. Crucially, the speaker probe gates of the K=2 fused logit pass cleanly: the literal 3-d fusion vector top-1 is $0.0182 \pm 0.0006$ and the 4167-d backbone-concat substrate is $0.0729 \pm 0.0005$, both well under the $0.0780$ ceiling.

\textbf{C2: A categorical closure of three de-confounding intervention classes on URTIC + frozen WavLM.} Data-level cross-speaker embedding mixup, after pivoting from a failed audio-splicing variant, gives a small UAR lift ($+0.006$) without measurably moving the speaker probe at conservative $\alpha$ and damages cold UAR at aggressive $\alpha$ (the operating window is empty across the $\alpha$-axis). Representation-level head-only supervised contrastive with speaker-masked positives is shown by random-projection and cold-CE controls to have no real mechanism beyond the bottleneck, and a combined-loss $\lambda$-sweep over $\{0, 0.05, 0.1, 0.25, 0.5\}$ shows pure cold-CE is monotonically Pareto-best on every metric. Gradient-level DANN/MDD adversary at both 128-d projection and 4096-d un-bottlenecked substrates is shown by per-epoch discriminator-health and frozen-substrate disc-ceiling diagnostics to have no recoverable speaker signal to act on; the bottleneck does the work the adversary was meant to add. None of the three classes beats the simpler cold-CE pipeline.

\textbf{C3: A methodology framework of seven negative-control disciplines.} Each discipline arose from a specific failure-mode the project hit and recovered from. M8: a self-splice control disambiguates ``cross-speaker mixing'' from ``the splicing operation itself'' on foundation-model features (URTIC's audio-level splice was detectable at $>$99\% UAR even within the same speaker). M9: $\alpha$-endpoint controls test both ends of an embedding-mixup design before declaring the rung locked (URTIC's mixup operating window is empty, not narrow). M10: the dimensionality bottleneck is a confound for any speaker-probe-based de-confounding measurement (a random-projection control at the same target dim is mandatory). M11: an auxiliary contrastive loss is purely subtractive on a CE-anchored bottleneck, not additive (verified by monotonic $\lambda$-sweep degradation). M12: a discriminator's train/devel memorisation gap is the in-flight validity check for adversarial training (high gap = pushing against memorised fingerprints, no useful signal). M13: the disc-ceiling diagnostic must run \emph{before} the adversarial $\lambda$-sweep, with $\textit{devel} > \textsc{HIGH}$ checked before any memo-gap-as-kill condition (the v1$\rightarrow$v2 worked example in this paper is a direct demonstration). M14: speaker-probe-as-de-confounding-measurement has a fundamental noise floor on small-data corpora with coarse pseudo-speaker labels — discriminator memorisation dominates regardless of bottleneck; report the disc-ceiling number alongside the probe number, and treat probe top-1 as an upper bound rather than a precise measurement. Each discipline generalises beyond URTIC.

\paragraph{Roadmap.} Section~\ref{sec:related} situates the work in the paralinguistic foundation-model and de-confounding literature. Section~\ref{sec:bg} covers URTIC, the frozen-WavLM cache, the speaker-grouped split discipline, and the speaker probe substrates. Section~\ref{sec:method} walks the controlled-positive pipeline (\textsc{A2} $\rightarrow$ \textsc{A2.5} $\rightarrow$ \textsc{A5b}). Section~\ref{sec:deconf} presents the de-confounding ladder closure with mechanism analyses. Section~\ref{sec:results} consolidates the headline numbers and recall-pattern shifts. Section~\ref{sec:disc} discusses limitations, the devel-vs-hidden-test caveat, and the exploratory-vs-confirmatory framing. Section~\ref{sec:conc} closes with the methodology argument.
